\documentclass[final, appendixprefix=true, chapterprefix=true, fontsize=12pt, numbers=noendperiod]{scrbook}
\newcommand{\myversion}{0.0}
\newcommand{\mytitle}{An introduction to Lippmann--Schwinger solvers}
\newcommand{\myauthor}{Sébastien Brisard}
\author{\myauthor}
\title{\mytitle}
\subtitle{Version~\myversion}

\addtokomafont{disposition}{\rmfamily}
\addtokomafont{caption}{\footnotesize}
\setkomafont{descriptionlabel}{\normalfont\bfseries}
\setcapindent{0em}

\usepackage{amsfonts}
\usepackage[fleqn]{amsmath}
\setlength{\mathindent}{1ex}
\usepackage{amssymb}

\usepackage{amsthm}
\theoremstyle{definition}
\renewcommand{\qedsymbol}{}
\newtheorem{remark}{Remark}[chapter]

\usepackage{comment}

\usepackage{graphicx}

\usepackage{polyglossia}
\setdefaultlanguage{english}

\usepackage[backend=biber,bibencoding=utf8,doi=false,giveninits=true,isbn=false,maxnames=10,minnames=5,sortcites=true,style=alphabetic-verb,texencoding=utf8,url=false]{biblatex}
\addbibresource{LS-intro.bib}

\usepackage{unicode-math}
\setmainfont{XITS}
\setmathfont{XITS Math}

\usepackage{xcolor}

\usepackage[allcolors=blue,breaklinks=true,colorlinks=true,linktocpage=true,pdfauthor={\myauthor},pdfsubject={},pdftitle={\mytitle}]{hyperref}

\usepackage{pdfcomment}
\defineavatar{SB}{author=SB,color=orange}

\usepackage{pict2e}

\newcommand{\cellindices}{\mathcal{P}}
\newcommand{\code}{\texttt}
\newcommand{\conj}{\operatorname{conj}}
\newcommand{\D}{\mathrm d}
\newcommand{\dbldot}{\mathbin{\mathord{:}}}
\newcommand{\dft}{\operatorname{DFT}}
\newcommand{\element}{\mathrm{e}}
\newcommand{\fftfreq}{\operatorname{Z}}
\newcommand{\I}{\mathrm i}
\newcommand{\integers}{\mathbb Z}
\newcommand{\naturals}{\mathbb N}
\newcommand{\reals}{\mathbb R}
\newcommand{\sinc}{\operatorname{sinc}}
\newcommand{\symotimes}{\stackrel{\mathrm{s}}{\otimes}}
\newcommand{\strain}{\boldsymbol{\varepsilon}}
\newcommand{\strains}{\mathcal E}
\newcommand{\stress}{\boldsymbol{\sigma}}
\newcommand{\stresses}{\mathcal S}
\newcommand{\tens}{\boldsymbol}
\newcommand{\tensors}{\mathcal T}
\newcommand{\tr}{\operatorname{tr}}
\newcommand{\tuple}{\mathsf}
\renewcommand{\vec}{\boldsymbol}

\begin{document}
\maketitle

\begin{center}
  \includegraphics{img/cc-by-nc-88x31.png}

  This work is licensed under a
  \href{https://creativecommons.org/licenses/by-nc/4.0/}{Creative Commons
    Attribution-NonCommercial 4.0 International License}.
\end{center}

\chapter{Nomenclature}

\begin{itemize}
\item \(d\): dimension of the physical space (typically \(d=2, 3\))
\item \(\Omega\): \(d\)-dimensional unit-cell
\item \(L_1,\ldots, L_d\): dimensions of the unit-cell:
  \(\Omega=(0, L_1)\times(0, L_2)\times\cdots\times(0, L_d)\)
\item \(\lvert\Omega\rvert=L_1L_2\cdots L_d\): volume of the unit-cell
\item \(\tuple{n}\): \(d\)-dimensional tuple of integers
  \(\tuple{n}=(n_1, n_2, \ldots, n_d)\)
\item \(\tuple{N}=(N_1, N_2, \ldots, N_d)\): size of the simulation grid
\item \(\lvert N\rvert=N_1N_2\cdots N_d\): total number of cells
\item \(h_i=L_i/N_i\): size of the cells (\(i=1, \ldots, d\))
\item
  \(\cellindices=\{0, \ldots, N_1-1\}\times\{0, \ldots, N_2-1\}\times\{0,
  \ldots, N_3-1\}\): set of cell indices
\item \(\Omega_{\tuple{p}}\): cells of the simulation grid
  (\(\tuple{p}\in\cellindices\))
\end{itemize}

%%% Local Variables:
%%% coding: utf-8
%%% fill-column: 80
%%% mode: latex
%%% TeX-engine: xetex
%%% TeX-master: "LS-intro.tex"
%%% End:


\chapter{Fourier transforms in a periodic setting}

This chapter provides a brief overview of the two Fourier transforms that are
going to be used in the previous document, namely Fourier series (see
Sec.~\ref{sec:20210924084626}) and discrete Fourier transforms (see
Sec.~\ref{sec:20210924084718}).

\section{Fourier series}
\label{sec:20210924084626}

Owing to the periodic setting, the fields that are involved in the various BVPs
to be discussed in this document are expanded in Fourier series. \(\tens T\)
being a \(\Omega\)-periodic tensor field (with sufficient regularity), the
following decomposition holds
\begin{equation}
  \tens T(\vec x)
  =\sum_{\tuple{n}\in\integers^d}\mathcal F(\tens T)(\vec k_{\tuple{n}})
  \exp(\I\vec k_{\tuple{n}}\cdot\vec x),
\end{equation}
where \(\tuple{n}\) denotes a \(d\)-dimensional tuple of integers (see
nomenclature, Chap.~\ref{cha:20210924084856}). The wave vectors
\(\vec k_{\tuple{n}}\) are given by
\begin{equation}
  \vec k_{\tuple{n}}=\frac{2\pi n_1}{L_1}\vec e_1
  +\frac{2\pi n_2}{L_2}\vec e_2
  +\cdots+\frac{2\pi n_d}{L_d}\vec e_d,
\end{equation}
and the Fourier coefficients of \(\tens T\) are defined as follows
\begin{equation}
\mathcal F(\tens T)(\vec k)=\frac1{\lvert\Omega\rvert}\int_{\vec x\in\Omega}
\tens T(\vec x)\exp(-\I\vec k\cdot\vec x)\,\D x_1\cdots\D x_d.
\end{equation}

It is recalled that the Fourier coefficients of the gradient and divergence of
\(\tens T\) can readily be computed from the Fourier coefficients of \(\tens T\)
\begin{equation}
\mathcal F(\tens T\otimes\nabla)(\vec k)=\mathcal F(\tens T)(\vec k)\otimes\I\vec k
\quad\text{and}\quad
\mathcal F(\tens T\cdot\nabla)(\vec k)=\mathcal F(\tens T)(\vec k)\cdot\I\vec k.
\end{equation}

When no confusion is possible, we will use a tilde to denote the Fourier
coefficients:
\(\tilde{\tens T}_\tuple{n}=\mathcal F(\tens T)(\vec k_\tuple{n})\).

\section{Discrete Fourier transforms}
\label{sec:20210924084718}

We consider a \(d\)-dimensional grid of size \(N_1\times\cdots\times N_d\). The
set of cell indices over this grid is
\begin{equation}
\cellindices=\{1,\ldots, N_1\}\times\cdots\times\{1, \ldots, N_d\},
\end{equation}
(refer to Remark~\ref{rem:20210924090334} for naming conventions of
indices). For \(\tuple{n}, \tuple{p}\in\cellindices\), we define
\(\Phi_{\tuple{n}\tuple{p}}\)
\begin{equation}
\phi_{\tuple{n}\tuple{p}}=2\pi\frac{\bigl(n_1-1\bigr)\bigl(p_1-1\bigr)}{N_1}
+\cdots+2\pi\frac{\bigl(n_d-1\bigr)\bigl(p_d-1\bigr)}{N_d}.
\end{equation}

Note that cell-indices start at 1, hence the \(-1\) correction in the above
formula (see Remark~\ref{rem:20210924090344}).

Let \(x=(x_{\tuple{p}})\) be a finite set of scalar values indexed by the
\(d\)-tuple \(\tuple{p}\in\cellindices\). The discrete Fourier transform is a
discrete set of scalar values \(\dft_{\tuple{n}}(x)\) indexed by the \(d\)-tuple
\(\tuple{n}\in\integers^d\), defined as follows
\begin{equation}
\dft_{\tuple{n}}(x)=\sum_{\tuple{p}\in\cellindices}
\exp\bigl(-\I\phi_{\tuple{n}\tuple{p}}\bigr)x_{\tuple{p}}.
\end{equation}

Note that in the above definition, no restrictions are applied to the
multi-index \(\tuple{n}\). However, it can be verified that the above series of
tensors is in fact \(\tuple{N}\)-periodic:
\(\dft_{\tuple{n}+\tuple{N}}(x)=\dft_{\tuple{n}}(x)\), where
\(\tuple{n}+\tuple{N}=(n_1+N_1, \ldots, n_d+N_d)\). Therefore, the
\(\tuple{n}\)-index is effectively restricted to \(1\leq n_i\leq N_i\) as
well. The most important results concerning the DFT are the \emph{inversion
  formula}
\begin{equation}
x_{\tuple{p}}=\frac1{\lvert\tuple{N}\rvert}\sum_{\tuple{n}\in\cellindices}
\exp\bigl(\I\phi_{\tuple{n}\tuple{p}}\bigr)\dft_{\tuple{n}}(x),
\end{equation}
the \emph{Plancherel theorem}
\begin{equation}
\sum_{\tuple{p}\in\cellindices}\conj(x_{\tuple p})y_{\tuple p}
=\frac1{\lvert\tuple N\rvert}\sum_{\tuple{n}\in\cellindices}
\conj[\dft_{\tuple n}(x)]\dft_{\tuple n}(y),
\end{equation}
and the \emph{circular convolution theorem}
\begin{equation}
\dft_{\tuple{n}}(x\ast y)=\dft_{\tuple{n}}(x)\dft_{\tuple{n}}(y),
\quad\text{where}\quad
(x\ast y)_{\tuple p}=\sum_{\tuple{q}\in\cellindices}
x_{\tuple{q}}y_{\tuple{p}-\tuple{q}+1}.
\end{equation}

The DFT is readily extended to tensor data points. In the absence of ambiguity,
the shorthand \(\hat{x}_{\tuple{n}}\) will be adopted for
\(\dft_{\tuple{n}}(x)\).

To close this section, we observe that the DFT of a series of \emph{real} data
points is a series of \emph{complex} data points. However, these complex values
have the following property
\begin{equation}
\dft_{\tuple{N}-\tuple{n}}(x)=\conj[\dft_{\tuple{n}}(x)],
\end{equation}
which is actually a \emph{necessary and sufficient} condition for the
\(x_{\tuple{p}}\) to be real.

%%% Local Variables:
%%% coding: utf-8
%%% fill-column: 80
%%% mode: latex
%%% TeX-engine: xetex
%%% TeX-master: "LS-intro.tex"
%%% End:


\chapter{Continuous Green operators}

In this chapter, we discuss various boundary-value problems in a periodic
setting. For each of these problems, we introduce the associated *Green
operator*.

\section{Conductivity}

\section{Continuous Green operator for linear elasticity}

We first define a few functional spaces; \(\tensors_2(\Omega)\) denotes the
space of second-order, symmetric, tensor fields, with square-integrable
components. Then, the space \(\stresses(\Omega)\) of periodic, self-equilibrated
stresses is defined as follows
\begin{equation}
  \tens\sigma\in\stresses(\Omega)\iff\left\{
    \begin{gathered}
      \tens\sigma\in\tensors_2(\Omega)\\
      \tens\sigma\cdot\nabla=\vec 0\text{ a.e in }\Omega\\
      \tens\sigma\cdot\vec e_i\text{ is }L_i\vec e_i
      \text{-periodic for all }i=1, 2, \ldots, d\text{ (no summation),}
    \end{gathered}
  \right.
\end{equation}
where the last condition expresses the periodicity of tractions in all
directions parallel to the sides of the unit-cell. The space
\(\tens\strains(\Omega)\) of periodic, geometrically compatible strains is
defined as follows
\begin{equation}
  \tens\varepsilon\in\strains(\Omega)\iff\left\{
    \begin{gathered}
      \tens\varepsilon\in\tensors_2(\Omega)\\
      \tens\varepsilon=\vec u\symotimes\nabla
      \text{ a.e. in }\Omega\text{ for some vector field }\vec u\\
      \vec u\text{ has square-integrable components}\\
      \vec u\text{ is }\Omega\text{-periodic.}
    \end{gathered}
  \right.
\end{equation}

Finally, we define the spaces of stresses and strains with zero average
\begin{equation}
  \stresses_0(\Omega)=\bigl\{\tens\sigma\in\stresses(\Omega),
  \langle\tens\sigma\rangle=\tens0\bigr\}
  \quad\text{and}\quad
  \strains_0(\Omega)=\{\tens\varepsilon\in\strains(\Omega),
  \langle\tens\varepsilon\rangle=\tens0\bigr\}.
\end{equation}

We are now ready to define the periodic, fourth-order Green operator for strains
\(\tens\Gamma\). Let \(\tens C\) be the homogeneous elastic stiffness of the
body \(\Omega\)\footnote{In other words, \(\tens C\) is a constant, fourth-order
  tensor with major and minor symmetries; furthermore, \(\tens C\) is positive
  definite.}. Let \(\tens\tau\in\tensors_2(\Omega)\) be a prescribed tensor
field (\emph{stress-polarization}). We want to find the equilibrium state of the
body \(\Omega\), subjected to the eigenstress \(\tens\tau\) and periodic boundary
conditions. In other words, we want to find the solution to the following
problem
\begin{equation}
  \text{Find }\tens\sigma\in\stresses_0(\Omega)
  \text{ and }\tens\varepsilon\in\strains_0(\Omega)
  \text{ such that }\tens\sigma=\tens C\dbldot\tens\varepsilon+\tens\tau
  \text{ a.e. in }\Omega.
\end{equation}

Owing to the periodic boundary conditions, we use Fourier series expansions of
\(\tens\tau\), \(\tens\sigma\), \(\tens\varepsilon\) and \(\vec u\) (see
Sec.~\ref{sec:20210924084626})
\begin{equation}
  \begin{Bmatrix}
    \tens\tau(\vec x)\\
    \tens\sigma(\vec x)\\
    \tens\varepsilon(\vec x)\\
    \vec u(\vec x)
  \end{Bmatrix}
  =\sum_{\tuple{n}\in\integers^d}
  \begin{Bmatrix}
    \tilde{\tens\tau}_\tuple{n}\\
    \tilde{\tens\sigma}_\tuple{n}\\
    \tilde{\tens\varepsilon}_\tuple{n}\\
    \tilde{\vec u}_\tuple{n}
  \end{Bmatrix}
  \exp(\I \vec k_\tuple{n}\cdot\vec x).
\end{equation}

The Fourier modes \(\tilde{\tens\sigma}_\tuple{n}\),
\(\tilde{\tens\varepsilon}_\tuple{n}\) and \(\tilde{\vec u}_\tuple{n}\) solve
the following equations (respectively: equilibrium, geometric compatibility,
constitutive relation)
\begin{gather}
  \label{eq:20210927092532}
  \tilde{\tens\sigma}_\tuple{n}\cdot\vec k_\tuple{n}=\vec 0\\
  \label{eq:20210927092538}
  \tilde{\tens\varepsilon}_\tuple{n}
  =\I\,\tilde{\vec u}_\tuple{n}\symotimes\vec k_\tuple{n}\\
  \label{eq:20210927092543}
  \tilde{\tens\sigma}_\tuple{n}=\tens C\dbldot\tilde{\tens\varepsilon}_\tuple{n}
  +\tilde{\tens\tau}_\tuple{n}.
\end{gather}

Plugging the Eq.~\eqref{eq:20210927092543} into Eq.~\eqref{eq:20210927092538},
and recalling that \(\tens C\) has the minor symmetries, we find the following
expression of \(\tilde{\tens\sigma}\)
\begin{equation}
  \tilde{\tens\sigma}_\tuple{n}
  =\I\bigl(\tens C\cdot\vec k_\tuple{n}\bigr)\cdot\tilde{\vec u}_\tuple{n}
  +\tilde{\tens\tau}_\tuple{n}.
\end{equation}

The Cauchy stress tensor being symmetric, Eq.~\eqref{eq:20210927092532} also
reads \(\vec k_\tuple{n}\cdot\tilde{\tens{\sigma}}_\tuple{n}=\vec 0\) and
\begin{equation}
  \tilde{\vec u}_\tuple{n}
  =\I\bigl(\vec k_\tuple{n}\cdot\tens C\cdot\vec k_\tuple{n}\bigr)^{-1}
  \cdot\tilde{\tens\tau}_\tuple{n}\cdot\vec k_\tuple{n},
\end{equation}
which delivers the following expression for the Fourier modes of the strain
field
\begin{equation}
  \tilde{\tens\varepsilon}_\tuple{n}
  =-\bigl[\bigl(\vec k_\tuple{n}\cdot\tens C\cdot\vec k_\tuple{n}\bigr)^{-1}
  \cdot\tilde{\tens\tau}_\tuple{n}\cdot\vec k_\tuple{n}\bigr]
  \symotimes\vec k_\tuple{n}.
\end{equation}
The above relation defines a linear mapping between
\(\tilde{\tens\tau}_\tuple{n}\) and \(\tilde{\tens\varepsilon}_\tuple{n}\). For
each Fourier mode \(\tuple{n}\), we therefore introduce the fourth-order tensor
\(\tilde{\tens\Gamma}_\tuple{n}\) with major and minor symmetries, such that
\(\tilde{\tens\varepsilon}_\tuple{n}=-\tilde{\tens\Gamma}_\tuple{n}\dbldot{\tilde{\tens\tau}}_\tuple{n}\). Our
analysis shows that
\(\tilde{\tens\Gamma}_\tuple{n}=\hat{\tens\Gamma}(\vec k_\tuple{n})\) where, for
arbitrary wave-vector \(\vec k\), \(\hat{\tens\Gamma}(\vec k)\) is a
fourth-order tensor with major and minor symmetries, such that
\begin{equation}
  \hat{\tens\Gamma}(\vec k)\dbldot\tilde{\tens\tau}
  =\bigl[\bigl(\vec n\cdot\tens C\cdot\vec n\bigr)^{-1}
  \cdot\tilde{\tens\tau}\cdot\vec n\bigr]\symotimes\vec n,
\end{equation}
where \(\vec n=\vec k/\lVert\vec k\rVert\). The above equation defines
\(\hat{\tens\Gamma}(\vec k)\) by how it operates on second-order, symmetric
tensors. A closed-form expression of this tensor can be derived in the case of
an isotropic material, for which
\begin{equation}
  \tens C=\lambda\tens I_2\otimes\tens I_2+2\mu\tens I_4,
\end{equation}
where \(\tens I_2\) (resp. \(\tens I_4\)) is the second-order
(resp. fourth-order) identity tensor, and \(\lambda\), \(\mu\) are the Lamé
coefficients. Then
\begin{equation}
  \vec n\cdot\bigl(\tens I_2\otimes\tens I_2\bigr)\cdot\vec n=\vec n\otimes\vec n
\end{equation}
and (recalling that \(\lVert\vec n\rVert=1\))
\begin{equation}
  \begin{aligned}[b]
    \vec n\cdot\tens I_4\cdot\vec n
    &=\tfrac12 n_i\bigl(\delta_{ik}\delta_{jl}+\delta_{il}\delta_{jk}\bigr)n_l
    \vec e_j\otimes\vec e_k
    =\tfrac12\bigl(n_kn_j+n_in_i\delta_{jk}\bigr)\vec e_j\otimes\vec e_k\\
    &=\tfrac12\bigl[\vec n\otimes\vec n+\bigl(\vec n\cdot\vec n\bigr)\tens I_2\bigr]
    =\tfrac12\bigl(\vec n\otimes\vec n+\tens I_2\bigr)
    =\vec n\otimes\vec n+\tfrac12\bigl(\tens I_2-\vec n\otimes\vec n\bigr)
  \end{aligned}
\end{equation}
and finally, we find the following expression of the acoustic tensor
\begin{equation}
  \vec n\cdot\tens C\cdot\vec n
  =\bigl(\lambda+2\mu\bigr)\vec n\otimes\vec n
  +\mu\bigl(\tens I_2-\vec n\otimes\vec n\bigr)
  =2\mu\frac{1-\nu}{1-2\nu}\vec n\otimes\vec n
  +\mu\bigl(\tens I_2-\vec n\otimes\vec n\bigr),
\end{equation}
where \(\nu\) denotes the Poisson ratio. The above second-order tensor is easily
inverted, since \(\vec n\otimes\vec n\) and \(\tens I_2-\vec n\otimes\vec n\)
are two orthogonal projectors (in the sense of the ``\(\dbldot\)'' product)
\begin{equation}
  2\mu\bigl(\vec n\cdot\tens C\cdot\vec n\bigr)^{-1}
  =\frac{1-2\nu}{1-\nu}\vec n\otimes\vec n
  +2\bigl(\tens I_2-\vec n\otimes\vec n\bigr)
  =2\tens I_2-\frac{\vec n\otimes\vec n}{1-\nu},
\end{equation}
from which it results that
\begin{equation}
  2\mu\bigl(\vec n\cdot\tens C\cdot\vec n\bigr)^{-1}
  \cdot\tilde{\tens\tau}\cdot\vec n
  =2\tilde{\tens\tau}\cdot\vec n
  -\frac{\vec n\cdot\tilde{\tens\tau}\cdot\vec n}{1-\nu}\vec n
\end{equation}
and we finally get
\begin{equation}
  \label{eq:20210927093916}
  2\mu\hat{\tens \Gamma}(\vec k)\dbldot\tilde{\tens \tau}
  =\bigl(\tilde{\tens \tau}\cdot\vec n\bigr)\otimes\vec n
  +\vec n\otimes\bigl(\tilde{\tens \tau}\cdot\vec n\bigr)
  -\frac{\vec n\cdot\tilde{\tens \tau}\cdot\vec n}{1-\nu}\vec n\otimes\vec n.
\end{equation}

The components of the \(\hat{\tens \Gamma}\) tensor are then readily found
\begin{equation}
  \hat{\Gamma}_{ijkl}(\vec k)
  =\frac{\delta_{ik}n_jn_l+\delta_{il}n_jn_k+\delta_{jk}n_in_l+\delta_{jl}n_in_k}{4\mu}
  -\frac{n_in_jn_kn_l}{2\mu\bigl(1-\nu\bigr)},
\end{equation}
which coincide with classical expressions (see
e.g. \textcite{suqu1990}). Implementation of the above equation is cumbersome;
it is only used for testing purposes. In Scapin, only the \emph{matvec} product
is required, and Eq.~\eqref{eq:20210927093916} was implemented.

\section{Hyperelasticity}


\chapter{Discrete Green operators}

In this chapter, we introduce various discretizations of the Green operator; we
will adopt the vocabulary of linear elasticity, although the concepts apply to
all the various physical models presented in Chap.~\ref{cha:20210928051857}.

\section{The \(\fftfreq\) function}

For \(n, N\in\naturals\), \(0\leq n<N\), we introduce \(\fftfreq(n, N)\)
\begin{equation}
  \fftfreq(n, N)=
  \begin{cases}
    n & \text{if }2n<N,\\
    n-N & \text{otherwise.}
  \end{cases}
\end{equation}

For \(n<0\) or \(n\geq N\), \(\fftfreq(n, N)\) is defined by
\(N\)-periodicity. \(\fftfreq\) is very similar to the NumPy
\href{https://numpy.org/doc/1.18/reference/generated/numpy.fft.fftfreq.html#numpy.fft.fftfreq}{fftfreq}
function. We have the important result
\begin{equation}
  \fftfreq(N-n, N)=
  \begin{cases}
    \fftfreq(n) & \text{if }2n=N,\\
    -\fftfreq(n) & \text{otherwise.}
  \end{cases}
\end{equation}
\begin{proof}
  To prove this property, several cases must be considered.
  \begin{enumerate}
  \item If \(n=0\)
    \begin{equation}
      \fftfreq(N-n, N)=\fftfreq(N, N)=\fftfreq(0, N)=0=-\fftfreq(n, N)
    \end{equation}
  \item If \(N\) is even, \(N=2M\)
    \begin{enumerate}
    \item If \(0<n<M\)
      \begin{gather*}
        2n<N\quad\Rightarrow\quad \fftfreq(n, N)=n,\\
        \begin{aligned}
          M<N-n&\quad\Rightarrow\quad N<2\bigl(N-n\bigr)\\
          &\quad\Rightarrow\quad \fftfreq(N-n, N)=N-n-N=-n.
        \end{aligned}
      \end{gather*}
    \item If \(n=M\)
      \begin{gather*}
        2n=N\quad\Rightarrow\quad \fftfreq(n, N)=n-N=-M,\\
        2(N-n)=2M=N\quad\Rightarrow\quad \fftfreq(N-n, N)=-M.
      \end{gather*}
    \item If \(M<n<N\)
      \begin{gather*}
        N<2n\quad\Rightarrow\quad \fftfreq(n, N)=n-N,\\
        \begin{aligned}
          N-n<M&\quad\Rightarrow\quad 2\bigl(N-n\bigr)<N\\
          &\quad\Rightarrow\quad \fftfreq(N-n, N)=N-n.
        \end{aligned}
      \end{gather*}
    \end{enumerate}
  \item If \(N\) is odd, \(N=2M+1\)
    \begin{enumerate}
    \item If \(0<n\leq M\)
      \begin{gather*}
        2n\leq 2M<N\quad\Rightarrow\quad \fftfreq(n, N)=n,\\
        \begin{aligned}
          M+1\leq N-n&\quad\Rightarrow\quad N<2\bigl(N-n\bigr)\\
          &\quad\Rightarrow\quad \fftfreq(N-n, N)=N-n-N=-n.
        \end{aligned}
      \end{gather*}
    \item  If \(M+1\leq n<N\)
      \begin{gather*}
        N+1\leq 2n\quad\Rightarrow\quad \fftfreq(n, N)=n-N\\
        \begin{aligned}
          N-n\leq N-M-1=M&\quad\Rightarrow\quad 2\bigl(N-n\bigr)\leq 2M<N\\
          &\quad\Rightarrow\quad \fftfreq(N-n, N)=N-n.
        \end{aligned}
      \end{gather*}
    \end{enumerate}
  \end{enumerate}
  Which completes the proof.
\end{proof}

The \(\fftfreq\) function can be defined for \(d\)-tuples as well
\begin{equation}
  \tuple \fftfreq(\tuple n, \tuple N)
  =\bigl(\fftfreq(n_1, N_1), \ldots, \fftfreq(n_d, N_d)\bigr)
\end{equation}
and we have again
\begin{equation}
  \tuple \fftfreq(\tuple N-\tuple n, \tuple N)=-\tuple \fftfreq(\tuple n)
\end{equation}
if none of the \(n_i\) is such that \(2n_i=N_i\).

\section{The approximation space}
\label{sec:20210928054031}

In order to define a discrete Green operator, we need to introduce the
approximation space for the stress-polarizations. We will consider here
stress-polarizations that are constant over each cell of a regular grid of size
\(N_1\times\cdots\times N_d\). The cells of this grid are
\begin{equation}
  \Omega_{\tuple{p}}^{\tuple{h}}=\{p_ih_i\leq x_i<\bigl(p_i+1\bigr)h_i, i=1,\ldots, d\},
\end{equation}
where \(\tuple{p}=(p_1,\ldots,p_d)\in\cellindices\) denotes a \(d\)-tuple of
integers, such that \(0\leq p_i<N_i\), \(i=1,\ldots, d\) and \(h_i=L_i/N_i\) is
the cell-size in the \(i\)-th direction. Note that, in the above expression, no
summation is implied by the repeated index \(i\). The total number of cells is
\(\lvert N\rvert=N_1\cdots N_d\).

We consider discrete stress-polarizations \(\tens\tau^{\tuple{h}}\) that are
constant over each cell of the grid: \(\tens\tau_{\tuple{p}}^{\tuple{h}}\)
denotes the constant value of \(\tens\tau^{\tuple{h}}\) in cell
\(\Omega_{\tuple{p}}^{\tuple{h}}\). The \(\tuple{h}\) superscript reminds that
\(\tens\tau^{\tuple{h}}\) is a discrete approximation of the true
stress-polarization \(\tens\tau\). We will call this approximation subspace:
\(\tensors_2^{\tuple{h}}(\Omega)\).

The discrete Green operator is defined as the restriction to this approximation
space of the continuous Green operator, seen as a bilinear form, or an
approximation of it. In other words, we want to propose an approximation of the
quantity
\begin{equation}
  \langle\tens\varpi^{\tuple{h}}\dbldot\tens\Gamma(\tens\tau^{\tuple{h}})\rangle
  \simeq\langle\tens\varpi^{\tuple{h}}\dbldot\tens\Gamma^{\tuple{h}}(\tens
  \tau^{\tuple{h}})\rangle\quad\text{for all }\tens\tau^{\tuple{h}},
  \tens\varpi^{\tuple{h}}\in\tensors_2^{\tuple{h}}(\Omega),
\end{equation}
where \(\tens\Gamma^{\tuple{h}}\) is defined only over
\(\tensors_2(\Omega)\). \(\tens\Gamma^{\tuple{h}}\) can therefore be seen as a
linear mapping between the cell values \(\tens\tau_{\tuple{p}}^{\tuple{h}}\) of
\(\tens\tau^{\tuple{p}}\) and the cell values of
\(\tens\Gamma^{\tuple{h}}(\tens\tau^{\tuple{h}})\); \(\tens\Gamma^{\tuple{h}}\)
is therefore a *matrix*, and the above equation should be understood as
\begin{equation}
  \label{eq:20210928054147}
  \langle\tens\varpi^{\tuple{h}}\dbldot\tens\Gamma(\tens\tau^{\tuple{h}})\rangle
  \simeq\frac1{\lvert\tuple{N}\rvert}\sum_{\tuple{p}, \tuple{q}\in\cellindices}
  \tens\varpi_{\tuple{p}}^{\tuple{h}}\dbldot
  \tens\Gamma_{\tuple{p}\tuple{q}}^{\tuple{h}}\dbldot\tens
  \tau_{\tuple{q}}^{\tuple{h}}
  \quad\text{for all}\quad
  \tens\tau^{\tuple{h}},\tens\varpi^{\tuple{h}}\in\tensors_2^{\tuple{h}}(\Omega).
\end{equation}

The continuous Green operator is translation invariant, and this property will
of course be transferred to the “exact” discrete Green operator; we will in fact
require *all* dicretizations of the Green operator to have this property. In
other words,
\(\tens\Gamma_{\tuple{p}\tuple{q}}^{\tuple{h}}=\tens\Gamma_{\tuple p-\tuple
  q}^{\tuple{h}}\) and the above equation reads
\begin{equation}
  \langle\tens\varpi^{\tuple{h}}\dbldot\tens\Gamma(\tens\tau^{\tuple{h}})\rangle
  \simeq\frac1{\lvert\tuple{N}\rvert}\sum_{\tuple{p}, \tuple{q}\in\cellindices}
  \tens\varpi_{\tuple{p}}^{\tuple{h}}\dbldot
  \tens\Gamma_{\tuple p-\tuple q}^{\tuple{h}}\dbldot\tens
  \tau_{\tuple{q}}^{\tuple{h}}\quad\text{for all }\tens\tau^{\tuple{h}},\tens
  \varpi^{\tuple{h}}\in\tensors_2^{\tuple{h}}(\Omega).
\end{equation}

Note that \(\tens\Gamma_{\tuple{p}}^{\tuple{h}}\) is now indexed by \emph{one}
index only and its DFT can be introduced unambiguously
\begin{equation}
  \begin{aligned}[b]
    &\frac{1}{\lvert\tuple N\rvert}\sum_{\tuple{p}, \tuple{q}\in\cellindices}
    \tens\varpi_{\tuple{p}}^{\tuple{h}}\dbldot\tens\Gamma_{\tuple p-\tuple q}^{\tuple h}
    \dbldot\tens\tau_{\tuple{q}}^{\tuple{h}}\\
    ={}&\frac1{\lvert\tuple N\rvert^2}\sum_{\tuple p, \tuple q, \tuple n\in\cellindices}
    \exp\Bigl[2\I\pi\sum_{j=1}^d\frac{n_j}{N_j}\bigl(p_j-q_j\bigr)\Bigr]
    \tens\varpi_{\tuple{p}}^{\tuple{h}}\dbldot\hat{\tens\Gamma}_{\tuple n}^{\tuple h}
    \dbldot\tens\tau_{\tuple q}^{\tuple h}\\
    ={}&\frac{1}{\lvert\tuple N\rvert^2}
    \sum_{\tuple n\in\cellindices}\Bigl[\sum_{\tuple{p}\in\cellindices}
    \exp\Bigl(2\I\pi\sum_{j=1}^d\frac{n_jp_j}{N_j}\Bigr)
    \tens\varpi_{\tuple p}^{\tuple h}\Bigr]\dbldot\hat{\tens\Gamma}_{\tuple n}^{\tuple h}
    \dbldot\Bigl[\sum_{\tuple q\in\cellindices}\exp\Bigl(-2\I\pi\sum_{j=1}^d
    \frac{n_jq_j}{N_j}\tens\tau_{\tuple q}^{\tuple{h}}\Bigr)\Bigr].
  \end{aligned}
\end{equation}

Since \(\tens\varpi\) is real, we have
\(\conj(\tens\varpi_{\tuple p}^h)=\tens\varpi_{\tuple p}^h\) and the first sum
in square brackets reads
\begin{equation}
  \sum_{\tuple{p}\in\cellindices}\exp\Bigl(2\I\pi\sum_{j=1}^d\frac{n_jp_j}{N_j}\Bigr)
  \tens\varpi_{\tuple p}^{\tuple h}=\conj\Bigl[\sum_{\tuple{p}\in\cellindices}
  \exp\Bigl(-2\I\pi\sum_{j=1}^d\frac{n_jp_j}{N_j}\Bigr)
  \tens\varpi_{\tuple p}^{\tuple h}\Bigr]
  =\conj(\hat{\tens\varpi}_{\tuple n}^{\tuple h}),
\end{equation}
while the second sum in square brackets reduces to \(\hat{\tens\tau}_{\tuple
  n}^{\tuple h}\). Gathering the above results, we find
\begin{equation}
  \frac{1}{\lvert\tuple N\rvert}\sum_{\tuple p, \tuple q\in\cellindices}
  \tens\varpi_{\tuple p}^{\tuple h}\dbldot\tens\Gamma_{\tuple p-\tuple q}^{\tuple h}
  \dbldot\tens\tau_{\tuple q}^{\tuple h}=\frac1{\lvert\tuple N\rvert^2}
  \sum_{\tuple n\in\cellindices}\conj(\hat{\tens\varpi}_{\tuple n}^{\tuple h})
  \dbldot\hat{\tens\Gamma}_{\tuple n}^{\tuple h}\dbldot
  \hat{\tens\tau}_{\tuple n}^{\tuple h}.
\end{equation}

The above equation can be understood as follows. \(\tens\Gamma^{\tuple h}\) is a
linear operator that maps the cell-wise constant field \(\tens\tau^{\tuple h}\)
to the cell-wise constant field \(\tens\eta^{\tuple h}\), the cell-values of
which are given by their DFT
\begin{equation}
  \tens{\eta}_{\tuple p}^{\tuple h}=
  \dft^{-1}_{\tuple p}(\hat{\tens\eta}_{\bullet}^{\tuple h}),\quad\text{with}
  \quad\hat{\tens\eta}_{\tuple n}^{\tuple h}
  =\hat{\tens\Gamma}_{\tuple n}^{\tuple h}\dbldot
  \hat{\tens\tau}_{\tuple n}^{\tuple h}.
\end{equation}

Then, from the Plancherel theorem
\begin{equation}
  \frac1{\lvert\tuple N\rvert^2}
  \sum_{\tuple n\in\cellindices}\conj(\hat{\tens\varpi}_{\tuple n}^{\tuple h})
  \dbldot\hat{\tens\Gamma}_{\tuple n}^{\tuple h}\dbldot
  \hat{\tens\tau}_{\tuple n}^{\tuple h}=\frac1{\lvert\tuple N\rvert^2}
  \sum_{\tuple n\in\cellindices}\conj(\hat{\tens\varpi}_{\tuple n}^{\tuple h})
  \dbldot\hat{\tens\eta}_{\tuple n}^{\tuple h}=\frac1{\lvert\tuple N\rvert}
  \sum_{\tuple p\in\cellindices}\tens\varpi_{\tuple p}^{\tuple h}
  \dbldot\tens\eta_{\tuple n}^{\tuple h}
\end{equation}
and the last sum can be seen as the volume average
\(\langle\tens\varpi^{\tuple h}\dbldot\tens\eta^{\tuple h}\rangle\). Remembering
that this expression was proposed as an approximation of
\(\langle\tens\varpi^{\tuple h}\dbldot\tens\Gamma(\tens\tau^{\tuple
  h})\rangle\), we finally find
\begin{equation}
  \langle\tens\varpi^{\tuple h}\dbldot\tens\Gamma(\tens\tau^{\tuple h})\rangle
  \simeq\langle\tens\varpi^{\tuple h}\dbldot\tens\eta^{\tuple h}\rangle
  =\langle\tens\varpi^{\tuple h}\dbldot\tens\Gamma^{\tuple h}(\tens\tau^{\tuple h})\rangle
  \quad\text{for all }\tens\varpi^{\tuple h}\in\tensors_2^{\tuple h}(\Omega),
\end{equation}
from which we find
\begin{equation}
  \tens\Gamma(\tens\tau^{\tuple h})\simeq\tens\Gamma^{\tuple h}
  (\tens\tau^{\tuple h}).
\end{equation}

The discrete Green operator, which was first introduced as an approximation of
the continuous Green operator, seen as a bilinear form, can also be understood
as an approximation of the continuous Green operator, seen as a linear
mapping. This latter point of view will become extremely efficient when it comes
to discretizing the Lippmann–Schwinger equation.

It results from the above developments that an explicit expression of the
discrete Green operator as a (gigantic) matrix is never needed. Instead, the
matrix-vector product
\(\tens\tau^{\tuple h}\mapsto\tens\Gamma^{\tuple h}(\tens\tau^{\tuple h})\) is
implemented in a matrix-free fashion as follows
\begin{enumerate}
\item Given \(\tens\tau^{\tuple h}\in\tensors_2^{\tuple h}(\Omega)\), compute
  the discrete Fourier transform \(\hat{\tens\tau}_{\tuple n}^{\tuple h}\) of
  its
  cell-values:\(\hat{\tens\tau}_{\tuple n}^{\tuple h}=\dft_{\tuple
    n}(\tens\tau_\bullet^{\tuple h})\),
\item for each discrete frequency, compute
  \(\hat{\tens\eta}_{\tuple n}^{\tuple h}=\hat{\tens\Gamma}_{\tuple n}^{\tuple
    h}\dbldot\hat{\tens\tau}_{\tuple n}^{\tuple h}\),
\item compute the inverse discrete Fourier transform
  \(\tens\eta_{\tuple{p}}^{\tuple h}\) of
  \(\hat{\tens\eta}_{\tuple{n}}^{\tuple h}\),
\end{enumerate}
discrete Fourier transforms being computed in steps 1 and 3 by means of the FFT.

\begin{remark}[Condition for the discrete Green operator to map real fields onto real fields]
  \label{rem:20210928060531}
  In the remainder of this chapter, we propose various discretizations of the
  Green operator. Before we proceed, though, it should be emphasized that the
  discrete Green operator must map a \emph{real} field onto a \emph{real}
  field. In other words, we must have
  \(\hat{\tens\eta}_{\tuple N-\tuple n}^{\tuple h}=\conj(\hat{\tens\eta}_{\tuple
    n}^{\tuple h})\) for all \(\tuple n\). Since
  \(\hat{\tens\eta}_{\tuple n}^{\tuple h} =\hat{\tens\Gamma}_{\tuple n}^{\tuple
    h}\dbldot\hat{\tens\tau}_{\tuple n}^{\tuple h}\) and
  \(\hat{\tens\tau}_{\tuple n}\) already satisfies this condition (it is the DFT
  of a \emph{real} field), we find the following requirement.

  Any discrete operator that will be considered below must ensure that
  \begin{equation}
    \hat{\tens\Gamma}_{\tuple N-\tuple n}^{\tuple h}
    =\conj(\hat{\tens\Gamma}_{\tuple n}^{\tuple h})
    \quad\text{for all}\quad\tuple n\in\cellindices.
  \end{equation}
\end{remark}


\section{Discretizations of the Green operator}

\subsection{The finite element discretization}

As discussed in Sec.~\ref{sec:20210928054031} [see in particular
Eq.~\eqref{eq:20210928054147}], the discrete Green operator can be seen as an
operator that delivers the cell-averages of the strains induced by a cell-wise
constant eigenstress \(\tens\tau^\tuple{h}\in\tensors_2^\tuple{h}(\Omega)\). In
other words, let us consider the solution to the following problem (see also
Sec.~\ref{sec:20210928054233})
\begin{equation}
  \label{eq:20210914140834}
  \text{Find }\tens\sigma\in\stresses_0(\Omega)
  \text{ and }\tens\varepsilon\in\strains_0(\Omega)
  \text{ such that }\tens\sigma=\tens C\dbldot\tens\varepsilon+\tens\tau^\tuple{h}
  \text{ a.e. in }\Omega.
\end{equation}

Note that, compared to the general problem that defines the continuous Green
operator (see Sec.~\ref{sec:20210928054233}), the above problem differs by the
space where \(\tens\tau^\tuple{h}\) lives. An approximate solution to this
problem can be retrieved from a finite element discretization
\parencite{bris2017}\footnote{The preprint of this paper is freely available on
  the \href{https://hal-enpc.archives-ouvertes.fr/hal-01304603}{HAL archive}; it
  is not the final version, but the theoretical sections were unchanged through
  the review process.}, where the mesh coincides with the grid introduced
previously and each element is discretized with linear shape functions (see
Sec.~\ref{}). The resulting linear system can be solved efficiently by means of
a matrix-free approach that is outlined below.

\subsubsection{The modal strain-displacement vector}

Introducing the nodal displacement \(\vec u_\tuple{n}\) at node
\(\tuple{n}\in\cellindices\), we define the inerpolated displacement
\(\vec u^\tuple{h}(\vec x)\) and strain \(\tens\strain^\tuple{h}(\vec x)\), as
well as the cell-average of the latter,
\(\overline{\tens\strain}_\tuple{p}^\tuple{h}\)
\begin{equation}
  \tens\varepsilon^\tuple{h}(\vec x)=\bigl(\vec
  u^\tuple{h}\symotimes\nabla\bigr)(\vec x)
  \quad\text{and}\quad
  \overline{\tens\varepsilon}_\tuple{p}^\tuple{h}
  =\frac1{\lvert\tuple{h}\rvert}\int_{\Omega_\tuple{p}}\tens\strain^\tuple{h}(\vec x)
  \,\D x_1\cdots\D x_d,
\end{equation}
where \(|\tuple{h}|=h_1\cdots h_d\) is the cell volume. In the present periodic
setting, the DFT \(\hat{\overline{\tens\varepsilon}}_\tuple{n}^\tuple{h}\) of
\(\tens\varepsilon\) can be expressed as follows
\begin{equation}
  \label{eq:20210914144114}
  \hat{\overline{\tens\varepsilon}}_\tuple{n}^\tuple{h}
  =\hat{\vec b}_\tuple{n}^\tuple{h}\symotimes\hat{\vec u}_\tuple{n}^\tuple{h},
\end{equation}
where \(\hat{\vec b}_\tuple{n}^\tuple{h}\) is the so-called \emph{modal
  strain-displacement vector} (see Ref.~\cite{bris2017} for its expression).

The modal strain-displacement vector introduced above is related to the *nodal*
strain-displacement operator introduced in Sec. [Gradient and
strain-displacement operators](@ref) of Chap. [On the d-dimensional brick
element](@ref). Indeed, in components, Eq. \eqref{eq20210914144114} reads
\begin{equation}
  \hat{\overline{\varepsilon}}_{ij\tuple{n}}^\tuple{h}
  =\tfrac12\bigl(\hat{b}_{i\tuple{n}}^\tuple{h}\hat{u}_{j\tuple{n}}^\tuple{h}
  +\hat{b}_{j\tuple{n}}^\tuple{h}\hat{u}_{i\tuple{n}}^\tuple{h}\bigr),
\end{equation}
and, from the circular convolution theorem
\begin{equation}
  \overline{\varepsilon}_{ij\tuple{p}}^\tuple{h}
  =\frac12\sum_{\tuple{q}\in\cellindices}\bigl(
  b_{i,\tuple{p}-\tuple{q}+1}^\tuple{h}u_{j\tuple{q}}^\tuple{h}
  +b_{j,\tuple{p}-\tuple{q}+1}^\tuple{h}u_{i\tuple{q}}^\tuple{h}\bigr)
  =\sum_{k=1}^d\sum_{\tuple{q}\in\cellindices}B_{ij\tuple{p}k\tuple{q}}^\tuple{h}
  u_{k\tuple{q}}^\tuple{h},
\end{equation}
where
\begin{equation}
  B_{ij\tuple{p}k\tuple{q}}^\tuple{h}
  =\frac12\bigl(\delta_{jk}b_{i,\tuple{p}-\tuple{q}+1}^\tuple{h}
  +\delta_{ik}b_{j,\tuple{p}-\tuple{q}+1}^\tuple{h}\bigr)
\end{equation}
is the global strain-displacement operator. The above identity reads, in Fourier
space
\begin{equation}
  \hat{B}_{ij\tuple{n}k\tuple{1}}^\tuple{h}
  =\frac12\bigl(\delta_{jk}\hat{b}_{i\tuple{n}}^\tuple{h}
  +\delta_{ik}\hat{b}_{j\tuple{n}}^\tuple{h}\bigr),
\end{equation}
where \(\tuple{q}\) has been set to \((1, \ldots, 1)\).

\subsubsection{The modal stiffness matrix}

It is recalled that the strain energy \(U\) is defined as the following integral
over the whole unit-cell \(\Omega\) (\(\lambda\), \(\mu\): Lamé coefficients)
\begin{equation}
  U=\frac12\int_\Omega\bigl[\lambda\bigl(\tr\tens\varepsilon^\tuple{h}\bigr)^2
  +2\mu\,\tens\varepsilon^\tuple{h}\dbldot\tens\varepsilon^\tuple{h}\bigr]
  \D x_1\cdots\D x_d.
\end{equation}

The strain energy is a quadratic form of the nodal displacements. Owing to the
homogeneity of the material (the Lamé coefficients are constant over the
unit-cell) and the periodic boundary conditions, the strain energy takes a
simple expression in Fourier space
\begin{equation}
  \label{eq:20210914060318}
  U=\frac 1{2\lvert\tuple{N}\rvert}
  \sum_{\tuple{n}\in\cellindices}
  \conj\bigl(\hat{\vec u}_\tuple{n}^\tuple{h}\bigr)
  \cdot\hat{\tens K}_\tuple{n}^\tuple{h}
  \cdot\hat{\vec u}_\tuple{n}^\tuple{h},
\end{equation}
where \(\hat{\tens K}_\tuple{n}^\tuple{h}\) is the \emph{modal stiffness matrix}.

\begin{remark}[Scaling of the modal stiffness matrix]
  The modal stiffness matrix introduced above differs from the matrix initially
  introduced in Ref.~\cite{bris2017} by a factor \(\lvert\tuple{h}\rvert\). More
  precisely,
  \begin{equation}
    \hat{\tens K}_\tuple{n}^\tuple{h}
    \bigr\rvert_\text{Scapin}
    =\lvert\tuple{h}\rvert\hat{\tens K}_\tuple{n}^\tuple{h}
    \bigr\rvert_\text{Brisard (2017)}.
  \end{equation}

  Such rescaling makes the relation between modal and nodal stiffness operators
  more natural (the former is the DFT of the latter).
\end{remark}

The modal stiffness matrix introduced above is related to the \emph{nodal}
(global) stiffness operator introduced in Sec.~\ref{} of Chap.~\ref{}. Indeed,
plugging the definition of the discrete Fourier transform (see
Sec.~\ref{sec:20210924084718}) into Eq. \eqref{eq:20210914060318} delivers the
following expression
\begin{equation}
  \begin{aligned}[b]
    U={}&\frac1{2\lvert\tuple{N}\rvert}\sum_{\tuple{n}, \tuple{p}, \tuple{q}\in\cellindices}
    \conj\bigl[\vec u_\tuple{p}^\tuple{h}\exp\bigl(-\I\phi_\tuple{np}\bigr)\bigr]
    \cdot\hat{\tens K}_\tuple{n}^\tuple{h}
    \cdot\vec u_\tuple{q}^\tuple{h}\exp\bigl(-\I\phi_\tuple{nq}\bigr)\\
    ={}&\frac1{2\lvert\tuple{N}\rvert}\sum_{\tuple{n}, \tuple{p}, \tuple{q}\in\cellindices}
    \conj\bigl(\vec u_\tuple{p}^\tuple{h}\bigr)
    \cdot\hat{\tens K}_\tuple{n}^\tuple{h}
    \cdot\vec u_\tuple{q}^\tuple{h}
    \exp\bigl[\I\bigl(\phi_\tuple{np}-\phi_\tuple{nq}\bigr)\bigr]\\
    ={}&\frac1{2\lvert\tuple{N}\rvert}\sum_{\tuple{n}, \tuple{p}, \tuple{q}\in\cellindices}
    \conj\bigl(\vec u_\tuple{p}^\tuple{h}\bigr)
    \cdot\hat{\tens K}_\tuple{n}^\tuple{h}
    \cdot\vec u_\tuple{q}^\tuple{h}
    \exp\bigl[\I\bigl(\phi_{\tuple{n}, \tuple{p}-\tuple{q}+1}\bigr)\bigr]\\
    ={}&\frac12\sum_{\tuple{p},\tuple{q}}
    \conj\bigl(\vec u_\tuple{p}^\tuple{h}\bigr)\cdot\tens K_\tuple{pq}^\tuple{h}
    \cdot\vec u_\tuple{q}^\tuple{h},
  \end{aligned}
\end{equation}
where
\begin{equation}
  \tens{K}_{\tuple{pq}}^\tuple{h}
  =\frac1{\lvert\tuple{N}\rvert}\sum_{\tuple{n}}
  \hat{\tens K}_\tuple{n}^\tuple{h}
  \exp\bigl(\I\phi_{\tuple{n},\tuple{p}-\tuple{q}+1}\bigr)
  =\dft_{\tuple{p}-\tuple{q}+1}^{-1}\bigl(\hat{\tens K}_\tuple{\bullet}^\tuple{h}\bigr)
\end{equation}
is the \emph{nodal stiffness matrix}, which appears as a block-circulant
matrix. This expresses the fact that the problem under consideration is
translation-invariant, owing to the homogeneity of the material and the periodic
boundary conditions (see also Sec.~\ref{}).

In order to derive a FE-approximate solution to the problem described by
Eq.~\eqref{eq:20210914140834}, we need to account for the contribution
\(U^\text{eigen}\) of the eigenstress \(\tens\tau^\tuple{h}\) to the total
potential energy \(\Pi\). This contribution reads, in the real space,
\begin{equation}
  U^\text{eigen}=\int_\Omega\bigl[\tens\tau^\tuple{h}(\vec x)
  \dbldot\tens\varepsilon(\vec x)\bigr]\D x_1\cdots\D x_d
\end{equation}
which can be expressed in Fourier space
\begin{equation}
  U^\text{eigen}=\frac{\lvert\tuple{h}\rvert}{\lvert\tuple{N}\rvert}
  \sum_{\tuple{n}\in\cellindices}\conj\bigl(\hat{\vec u}_\tuple{n}\bigr)
  \cdot\hat{\tens\tau}_\tuple{n}^\tuple{h}\cdot
  \conj\hat{\tens B}_\tuple{n}^\tuple{h},
\end{equation}
where \(\hat{\tens B}_\tuple{n}^\tuple{h}\) is the modal strain-displacement
operator introduced above, and \(\hat{\tens\tau}_\tuple{n}^\tuple{h}\) is the
DFT of the cell-values of the cell-wise constant eigenstrain.

Optimization of \(\Pi\) w.r.t. the nodal displacements delivers the following
equations
\begin{equation}
  \label{eq:20210914145137}
  \hat{\tens K}_\tuple{n}^\tuple{h}\cdot\hat{\vec u}_\tuple{n}^\tuple{h}
  =-\hat{\tens\tau}_\tuple{n}^\tuple{h}\cdot\hat{\tens B}_\tuple{n}^\tuple{h},
\end{equation}

Note that these equations are \emph{uncoupled}: they reduce to a \(d\times d\)
linear system \emph{for each Fourier mode}. The solution to these equations
delivers the modal displacements
\begin{equation}
  \hat{\vec u}_\tuple{n}^\tuple{h}
  =-\bigl(\hat{\tens K}_\tuple{n}^\tuple{h}\bigr)^{-1}
  \cdot\hat{\tens\tau}_\tuple{n}^\tuple{h}
  \cdot\hat{\tens B}_\tuple{n}^\tuple{h}.
\end{equation}

Combining with Eq.~\eqref{eq:20210914144114}, we get the following expression of
the cell-average of the strains, in Fourier space
\begin{equation}
  \hat{\overline{\tens\varepsilon}}_\tuple{n}^\tuple{h}
  =-\hat{\vec B}_\tuple{n}^\tuple{h}\symotimes
  \bigl[\bigl(\hat{\tens K}_\tuple{n}^\tuple{h}\bigr)^{-1}
  \cdot\hat{\tens\tau}_\tuple{n}^\tuple{h}\cdot\hat{\vec B}_\tuple{n}^\tuple{h}\bigr].
\end{equation}

Upon appropriate symmetrization (\(\tens\tau^\tuple{h}\) and
\(\tens\varepsilon^\tuple{h}\) are both symmetric, second-rank tensors), we find
\begin{equation}
  \hat{\overline{\tens\varepsilon}}_\tuple{n}^\tuple{h}=-\hat{\tens\Gamma}_\tuple{n}^{\tuple{h}, \text{Bri17}}\dbldot\hat{\tens\tau}_\tuple{n}^\tuple{h},
\end{equation}
with
\begin{equation}
  \hat{\tens\Gamma}_\tuple{n}^{\tuple{h}, \text{Bri17}}
  =\hat{\vec B}_\tuple{n}^\tuple{h}
  \symotimes\bigl(\tens K_\tuple{n}^\tuple{h}\bigr)^{-1}
  \symotimes\hat{\vec B}_\tuple{n}^\tuple{h},
\end{equation}
which defines the discrete Green operator.

\begin{remark}
  Eq. \eqref{eq:20210914145137} is singular for \(\vec k = \vec 0\). Indeed, in a
  periodic setting, the displacement is defined up to a constant translation.  It
  is convenient to select the solution with zero average, that is
  \(\hat{\vec u}_\tuple{0}^\tuple{h}=\vec 0\).
\end{remark}

\subsection{The discretization of Brisard and Dormieux}

It was proved by Brisard and Dormieux [(2010)](@ref bris2010a) that, for all
\(\tens\tau^{\tuple h}, \tens\varpi^{\tuple h}\in\tensors_2^{\tuple h}(\Omega)\)
\begin{equation}
  \langle\tens\varpi^h\dbldot\tens\Gamma(\tens\tau^h)\rangle=\frac1{N^2}
  \sum_{n_1=0}^{N_1-1}\cdots\sum_{n_d=0}^{N_d-1}\conj(\hat{\tens\varpi}_
  {\tuple{n}}^h)\dbldot\hat{\tens\Gamma}_n^{h, \mathrm{BD10}}\dbldot
  \hat{\tens\tau}_{\tuple{n}}^h,
\end{equation}
where
\begin{equation}
  \hat{\tens\Gamma}_{\tuple n}^{\tuple h, \mathrm{BD10}}
  =\sum_{\tuple m\in\integers^d}\bigl[F(\vec\alpha_{\tuple n+\tuple m\tuple N})
  \bigr]^2\hat{\tens\Gamma}(\vec k_{\tuple n+\tuple m\tuple N}),
\end{equation}
where \(\tuple{n+mN}\) denotes the \(d\)-tuple:
\(\tuple n+\tuple m\tuple N=(n_1+m_1N_1, \ldots, n_d+m_dN_d)\), while
\(\alpha_{\tuple{n}}\) is the following dimensionless vector
\begin{equation}
  \vec\alpha_{\tuple{n}}
  =\frac{2\pi h_1n_1}{L_1}\vec e_1+\cdots+\frac{2\pi h_dn_d}{L_d}\vec e_d,
\end{equation}
finally, \(F\) is the tensor product of sine cardinal functions
\begin{equation}
  F(\vec\alpha)=\sinc\frac{\alpha_1}2\cdots\sinc\frac{\alpha_d}2.
\end{equation}

Note that the above equations deliver the \emph{exact} cell-averages of the
Green operator, applied exactly to any \emph{cell-wise constant} polarization
field. Unfortunately, the series that defines the discrete Green operator can in
general not be evaluated, owing to very slow convergence. Therefore, this
discrete Green operator is unpractical, and is recalled here only for pedagocial
reasons.

\subsection{The discretization of Moulinec and Suquet}

This is probably the most simple discretization, introduced first by Moulinec
and Suquet~\cite{moul1994, moul1998}. Only the lowest (positive and negative)
frequencies are kept
\begin{equation}
  \hat{\tens\Gamma}_{\tuple n}^{\tuple h, \mathrm{MS94}}=\hat{\tens\Gamma}
  (\vec k_{\tuple Z(\tuple n, \tuple N)}).
\end{equation}

We must check that the condition for the discrete Green operator to map real
fields onto real fields is satisfied (see
Remark~\ref{rem:20210928060531}). Using the properties of the \(\fftfreq\)
function (see Sec.~\ref{}) and assuming first that none of the \(n_i\) is such
that \(2n_i=N_i\)
\begin{equation}
  \hat{\tens\Gamma}_{\tuple N-\tuple n}^{\tuple h, \mathrm{MS94}}
  =\hat{\tens\Gamma}(\vec k_{\tuple Z(\tuple N-\tuple n, \tuple N)})
  =\hat{\tens\Gamma}(\vec k_{-\tuple Z(\tuple n, \tuple N)})
  =\hat{\tens\Gamma}(-\vec k_{\tuple Z(\tuple n, \tuple N)})
\end{equation}

All Green operators presented in Chap.~\ref{cha:20210928051857} are such that
\(\hat{\tens\Gamma}(-\vec k)=\hat{\tens\Gamma}(\vec k)\), therefore
\begin{equation}
  \hat{\tens\Gamma}_{\tuple N-\tuple n}^{\tuple h, \mathrm{MS94}}
  =\hat{\tens\Gamma}(\vec k_{\tuple Z(\tuple n, \tuple N)})
  =\hat{\tens\Gamma}_{\tuple n}^{\tuple h, \mathrm{MS94}}
\end{equation}
and the property is verified. Conversely, if all the \(n_i\) are such that
\(2n_i=N_i\), then
\begin{equation}
  \hat{\tens\Gamma}_{\tuple N-\tuple n}^{\tuple h, \mathrm{MS94}}
  =\hat{\tens\Gamma}(\vec k_{\tuple Z(\tuple N-\tuple n, \tuple N)})
  =\hat{\tens\Gamma}(\vec k_{\tuple Z(\tuple n, \tuple N)})
  =\hat{\tens\Gamma}_{\tuple n}^{\tuple h, \mathrm{MS94}}.
\end{equation}

More problematic is the case when a few, but not all, \(n_i\) are such that
\(2n_i=N_i\). Then the property does not hold for such frequencies. Moulinec and
Suquet~\cite{moul1998} use a specific treatment for such cases
\begin{equation}
  \hat{\tens\Gamma}(\vec k_{\tuple n})=\tens C^{-1},
\end{equation}
if one of the \(n_i\) is such that \(2n_i=N_i\). This is implemented in
\texttt{Scapin}. Note that such cases occur only for even-sized grids.


\chapter{On the d-dimensional brick element}
\label{cha:20210930054817}

In this chapter, we formulate the linear brick finite element for conductivity
and linear elasticity in \(d\) dimensions (\(d\in\{2, 3\}\)).

\begin{remark}
  In the remainder of this chapter, \(i\), \(j\), \(h\) and \(k\) denote scalar
  indices that refer to tensor components, and span \(\{1, 2, \cdots d\}\);
  \(\tuple{ p}\), \(\tuple{q}\) are \(d\)-dimensional multi-indices that refer
  to nodes. Depending on the context (local, element operators or global,
  assembled operators) they span \(\{1, 2\}^d\) or the whole grid
  \(\cellindices\). We adopt Einstein's convention for indices of both types.
\end{remark}

\section{Geometry of the reference brick element}

The reference element being centered at the origin, it occupies the following
domain
\begin{equation}
  \Omega^\element=\biggl(-\frac{h_1}2, \frac{h_1}2\biggr)\times\cdots
  \times\biggl(-\frac{h_d}2, \frac{h_d}2\biggr),
\end{equation}
where \(h_i\) is the size of the element along axis \(i\) (\(i=1, \dots,
d\)). The nodes of the element are indexed by \(\tuple{p}\in\{1, 2\}^d\), such
that the coordinates of node \(\tuple{p}\) read
\begin{equation}
  x_{i\tuple{p}}=(-1)^{p_i}\frac{h_i}2.
\end{equation}

For integration purposes, it is useful to introduce the \emph{reduced
  coordinates} \(\xi_i=2x_i/h_i\)
\begin{equation}
  \int_{\Omega^\element}f(x_1, \ldots, x_d)\,\D x_1\cdots\D x_d
  =\frac{h_1\cdots h_d}{2^d}\int_{(-1, 1)^d}f(h_1\xi_1/2, \ldots, h_d\xi_d/2)
  \,\D\xi_1\cdots\D\xi_d
\end{equation}
and it is observed that the reduced coordinates of the node \(\tuple{p}\) are
\(\xi_i=(-1)^{p_i}\).

\section{Shape functions}
\label{sec:20210930054720}

The shape function associated with node \(\tuple{p}\) is
\(N^\element_{\tuple{p}}(\vec x)\), wich is defined below as a function of the
reduced coordinates of the observation point
\begin{equation}
  N^\element_{\tuple{p}}(\vec x)=\frac1{2^d}\prod_{i=1}^d\bigl[1+(-1)^{p_i}\xi_i\bigr]
\end{equation}
and we have as expected
\(N^\element_{\tuple{p}}(\vec x_{\tuple q}) = \delta_{\tuple{p}\tuple{q}}\). The
nodal values of \(f\) being \(f_\tuple{p}\), we have the interpolation formula:
\(f(\vec x)=N^\element_\tuple{p}(\vec x)f_\tuple{p}\).

\section{Gradient and strain-displacement operators}
\label{sec:20210930054808}

We consider a scalar interpolated field,
\(f(\vec x)=N^\element_\tuple{p}(\vec x)f_\tuple{p}\). The components of its
gradient are given by the following expression
\begin{equation}
  \bigl(f\otimes\nabla\bigr)_i=\partial_i f=D_{i\tuple{p}}^\element f_\tuple{p},
\end{equation}
where the observation point \(\vec x\) has been dropped, and
\(D_{i\tuple{p}}^\element\) denotes the element gradient operator
\begin{equation}
  D_{i\tuple{p}}^\element=\partial_iN_\tuple{p}
  =\frac{(-1)^{p_i}}{2^{d-1}h_i}
  \prod_{j\neq i}\bigl[1+(-1)^{p_j}\xi_j\bigr].
\end{equation}

It will be useful to introduce the element average of this operator,
\(\overline{D}_{i\tuple{n}}^\element\)
\begin{equation}
  \overline{D}_{i\tuple{p}}^\element=\frac 1{h_1\cdots h_d}
  \int_{\Omega^\element}D^\element_{i\tuple{p}}(\vec x)\,\D x_1\cdots\D x_d
\end{equation}
and, observing that each factor in square brackets in the above expression of
\(D^\element_{i\tuple{p}}\) varies linearly between 0 et 2 (and therefore
averages to 1), we find
\begin{equation}
  \overline{D}^\element_{i\tuple{p}}=\frac{(-1)^{p_i}}{2^{d-1}h_i}.
\end{equation}

We now consider the interpolated displacement field,
\(\vec u = N^\element_\tuple{p} \vec u_\tuple{p}\). The components of the
associated strain, \(\tens\varepsilon=\vec u\symotimes\nabla\) are
\begin{equation}
  \varepsilon_{ij}=\tfrac12\bigl(\partial_iu_j+\partial_ju_i\bigr)
  =\tfrac12\bigl(D^\element_{i\tuple{p}}u_{j\tuple{p}}
  +D^\element_{j\tuple{p}}u_{i\tuple{p}}\bigr),
\end{equation}
where \(u_{i\tuple{p}}\) denotes the \(i\)-th component of the nodal value of
\(\vec u\) at node \(\tuple{p}\),
\(u_{i\tuple{p}}=\vec e_i\cdot\vec u(\vec x_\tuple{p})\).

\begin{remark}
  In the Julia implementation of the brick element, the nodal displacements are
  stored in a \(d+1\)-dimensional array \code{u}, such that
  \(\text{\code{u[i, p]}}=u_{i\tuple{p}}\), where \code{i~in~\{1, ... d\}} and
  \code{p~in~CartesianIndices(1:N[1], …, 1:N[d])}. With this convention, the
  component index, \code{i}, is the \emph{fast} index. Therefore, all dofs of
  the same node are contiguous, which is the usual convention in finite element.
\end{remark}

Rearranging the above expression, we find
\(\varepsilon_{ij}=B^\element_{ijk\tuple{p}}u_{k\tuple{p}}^\element\), where
\begin{equation}
  B_{ijk\tuple{p}}^\element
  =\tfrac12\bigl(D^\element_{i\tuple{p}}\delta_{jk}
  +D^\element_{j\tuple{p}}\delta_{ik}\bigr)
\end{equation}
is the element strain-displacement operator, with volume average
\begin{equation}
  \overline{B}_{ijk\tuple{p}}^\element
  =\tfrac12\bigl(\overline{D}^\element_{i\tuple{p}}\delta_{jk}
  +\overline{D}^\element_{j\tuple{p}}\delta_{ik}\bigr)
  =\frac1{2^d}\Bigl[\frac{(-1)^{p_i}\delta_{jk}}{h_i}
  +\frac{(-1)^{p_j}\delta_{ik}}{h_j}\Bigr].
\end{equation}

It is obverved that the \emph{trace} of the strain tensor is given by the
following expression
\begin{equation}
  \varepsilon_{ii}=D_{i\tuple{p}}u_{i\tuple{p}}^\element.
\end{equation}

\section{Element stiffness operator}

\subsection{Conductivity}

\subsection{Linear elasticity}
\label{sec:20210930054950}

Assuming the material to be linearly elastic, with Lamé coefficients \(\lambda\)
and \(\mu\), we have the following expression of the interpolated stresses as a
function of the nodal displacements \(u_{i\tuple{p}}\)
\begin{equation}
  \sigma_{hk}=\bigl(\lambda\delta_{hk} D^\element_{i\tuple{p}}
  +2\mu B^\element_{hki\tuple{p}}\bigr)u_{i\tuple{p}}
\end{equation}
and the volume density of elastic energy reads
\begin{equation}
  \tfrac12\sigma_{hk}\varepsilon_{hk}
  =\bigl(\lambda\delta_{hk} D^\element_{i\tuple{p}}
  +2\mu B^\element_{hki\tuple{p}}\bigr)B_{hkj\tuple{q}}^\element
  u_{i\tuple{p}}u_{j\tuple{q}}.
\end{equation}

Upon expansion and integration, we get the expression of the elastic energy
within the element
\begin{equation}
  U^\element
  =\int_{\Omega^\element}\tfrac12\sigma_{hk}\varepsilon_{hk}\,\D x_1\cdots\D x_d
  =\tfrac12K^\element_{i\tuple{p}j\tuple{q}}u_{i\tuple{p}}u_{j\tuple{q}},
\end{equation}
where the element stiffness operator reads
\begin{equation}
  K^\element_{i\tuple{p}j\tuple{q}}
  =\lambda\,K^{\element, \lambda}_{i\tuple{p}j\tuple{q}}
  +\mu\,K^{\element, \mu}_{i\tuple{p}j\tuple{q}},
\end{equation}
with
\begin{equation}
  K^{\element, \lambda}_{i\tuple{p}j\tuple{q}}
  =\int_{\Omega^\element}D^\element_{i\tuple{p}}D^\element_{j\tuple{q}}
  \quad\text{and}\quad
  K^{\element, \mu}_{i\tuple{p}j\tuple{q}}
  =\int_{\Omega^\element} 2 B^\element_{hki\tuple{p}}B^\element_{hkj\tuple{q}}.
\end{equation}


\printbibliography

\end{document}

%%% Local Variables:
%%% coding: utf-8
%%% fill-column: 80
%%% mode: latex
%%% TeX-engine: xetex
%%% TeX-master: t
%%% End:
